\documentclass[11pt]{article}
\usepackage[margin=1in]{geometry}
\usepackage{amsmath}
\usepackage{amssymb}
\usepackage{enumitem}
\usepackage{xcolor}
\usepackage{tcolorbox}

\definecolor{boxblue}{RGB}{230,240,255}
\definecolor{boxgreen}{RGB}{230,255,240}

\title{\textbf{What are MC and AC?}}
\author{ECO310: Basic Cost Concepts}
\date{}

\begin{document}
\maketitle

\section{MC = Marginal Cost}
\textit{(Lecture 11-12, slide 7)}

\begin{tcolorbox}[colback=boxblue, colframe=blue!50!black, title=Definition:]
\Large
\[
MC(q) = C'(q) = \frac{dC}{dq}
\]

\normalsize
\textbf{Marginal Cost} = the derivative of the cost function

= the cost of producing \textbf{one more unit}
\end{tcolorbox}

\subsection*{What does it mean?}

MC tells you: \textbf{``How much does it cost to make one more unit?''}

\subsection*{Example:}

If $C(q) = 2q^2 + 10$, then:
\[
MC(q) = \frac{dC}{dq} = 4q
\]

\textbf{Interpretation:}
\begin{itemize}[leftmargin=*]
    \item When you've made $q=0$ units, making the 1st unit costs $MC(0) = 0$
    \item When you've made $q=5$ units, making the 6th unit costs $MC(5) = 20$
    \item When you've made $q=10$ units, making the 11th unit costs $MC(10) = 40$
\end{itemize}

Notice: It gets more expensive to make each additional unit! This is common in economics.

\newpage

\section{AC = Average Cost}
\textit{(Lecture 11-12, slide 28)}

\begin{tcolorbox}[colback=boxgreen, colframe=green!50!black, title=Definition:]
\Large
\[
AC(q) = \frac{C(q)}{q}
\]

\normalsize
\textbf{Average Cost} = total cost divided by quantity

= the cost \textbf{per unit} when making $q$ units
\end{tcolorbox}

\subsection*{What does it mean?}

AC tells you: \textbf{``On average, how much does each unit cost to make?''}

\subsection*{Example:}

If $C(q) = 2q^2 + 10$, then:
\[
AC(q) = \frac{C(q)}{q} = \frac{2q^2 + 10}{q} = 2q + \frac{10}{q}
\]

\textbf{Interpretation:}
\begin{itemize}[leftmargin=*]
    \item To make 1 unit total: $AC(1) = 2(1) + \frac{10}{1} = 12$ per unit
    \item To make 5 units total: $AC(5) = 2(5) + \frac{10}{5} = 12$ per unit
    \item To make 10 units total: $AC(10) = 2(10) + \frac{10}{10} = 21$ per unit
\end{itemize}

Notice: Average cost changes as you make more units!

\section{The Difference: MC vs. AC}

\begin{center}
\begin{tabular}{p{7cm}|p{7cm}}
\hline
\textbf{Marginal Cost (MC)} & \textbf{Average Cost (AC)} \\
\hline
Cost of \textbf{one more} unit & Cost \textbf{per unit} on average \\[0.3cm]
$MC = \frac{dC}{dq}$ (derivative) & $AC = \frac{C(q)}{q}$ (total/quantity) \\[0.3cm]
``How much to make the next one?'' & ``How much does each one cost?'' \\
\hline
\end{tabular}
\end{center}

\subsection*{Real-world analogy:}

Imagine a bakery making cookies:

\begin{itemize}[leftmargin=*]
    \item \textbf{Fixed costs:} \$100 (oven, rent, etc.)
    \item \textbf{Variable costs:} \$2 per cookie (flour, sugar, etc.)
\end{itemize}

So $C(q) = 100 + 2q$

\textbf{Marginal Cost:}
\[
MC(q) = \frac{dC}{dq} = 2
\]
``Each additional cookie costs \$2 to make''

\textbf{Average Cost:}
\[
AC(q) = \frac{C(q)}{q} = \frac{100 + 2q}{q} = \frac{100}{q} + 2
\]

\begin{itemize}[leftmargin=*]
    \item Make 1 cookie: $AC(1) = \frac{100}{1} + 2 = \$102$ per cookie (ouch!)
    \item Make 10 cookies: $AC(10) = \frac{100}{10} + 2 = \$12$ per cookie
    \item Make 100 cookies: $AC(100) = \frac{100}{100} + 2 = \$3$ per cookie
\end{itemize}

The fixed cost (\$100) gets spread over more cookies, so AC goes down!

\newpage

\section{Why Do We Care?}

\subsection*{For Problem Set 3:}

\textbf{Short-Run Equilibrium (Question 2a):}
\begin{itemize}[leftmargin=*]
    \item Firms produce where \textbf{p = MC}
    \item This maximizes profit for each firm
\end{itemize}

\textbf{Long-Run Equilibrium (Question 2b):}
\begin{itemize}[leftmargin=*]
    \item Firms still produce where \textbf{p = MC} (profit max)
    \item But also: \textbf{p = min AC} (zero profit condition)
    \item So we need: \textbf{MC = AC} to find equilibrium!
\end{itemize}

\section{Your Problem Set Example}

For Problem Set 3, Question 2:

\textbf{Given:} $C(q) = f(N) + \alpha q^2$

\subsection*{Find MC:}
\begin{align*}
MC(q) &= \frac{dC}{dq} = \frac{d}{dq}\left(f(N) + \alpha q^2\right) \\
&= 0 + 2\alpha q \\
&= 2\alpha q
\end{align*}

\subsection*{Find AC:}
\begin{align*}
AC(q) &= \frac{C(q)}{q} = \frac{f(N) + \alpha q^2}{q} \\
&= \frac{f(N)}{q} + \alpha q
\end{align*}

\subsection*{For Long-Run Equilibrium, set MC = AC:}
\begin{align*}
2\alpha q &= \frac{f(N)}{q} + \alpha q \\
\alpha q &= \frac{f(N)}{q} \\
\alpha q^2 &= f(N) \\
q^* &= \sqrt{\frac{f(N)}{\alpha}}
\end{align*}

That's how you solve it!

\section{Quick Reference}

\begin{tcolorbox}[colback=yellow!10, colframe=orange!75!black, title=Summary:]
\begin{align*}
\text{Cost Function: } & C(q) = \text{total cost of producing } q \text{ units} \\[0.3cm]
\text{Marginal Cost: } & MC(q) = C'(q) = \text{cost of one more unit} \\[0.3cm]
\text{Average Cost: } & AC(q) = \frac{C(q)}{q} = \text{cost per unit} \\[0.5cm]
\text{Profit Max: } & p = MC \\[0.3cm]
\text{Zero Profit (LR): } & p = \min AC \\[0.3cm]
\text{LR Equilibrium: } & MC = AC
\end{align*}
\end{tcolorbox}

\end{document}
